\documentclass[10pt,a4paper]{amsart}
\usepackage[margin=19mm]{geometry}
\usepackage{amsmath,amssymb,mathtools}
\usepackage[scr=rsfs,cal=euler]{mathalfa}
\usepackage{xfrac}
\usepackage[unicode,hidelinks,bookmarks=false]{hyperref}
\newcommand{\ie}{i.e.\ }
\newcommand{\cf}{cf.\ }
\newcommand{\resp}{respectively\ }
\newcommand{\Subsection}{\subsection}
\providecommand{\nobreakdash}{\nobreak\hbox{-}}
\setlength{\emergencystretch}{2em}
\multlinegap0pt
\usepackage{bm}
\usepackage{bbm}
\usepackage{dsfont}
\usepackage{xspace}
\usepackage{cancel}
\usepackage{mathtools}
\usepackage{amsthm}
\makeatletter
\@ifundefined{thm}{\newtheorem{thm}{Theorem}}{}
\makeatother
\DeclareMathOperator{\ehr}{ehr}
\DeclareMathOperator{\vol}{vol}
\newcommand{\Q}{\mathbb{Q}}
\newcommand{\R}{\mathbb{R}}
\newcommand{\Z}{\mathbb{Z}}
\renewcommand{\phi}{\varphi}
\title[Ehrhart Theory over Abelian Group Rings]{Ehrhart Theory over Abelian Group Rings}
\author{Robert Davis, Jes\'us A. De Loera, Alexey Garber, Katharina Jochemko, Mohamed Omar, Josephine Yu}
\date{Source: arXiv:2511.10373v2 (CC BY 4.0)}
\begin{document}
\maketitle

\noindent\emph{Source and licence.} Ehrhart Theory over Abelian Group Rings. Version arXiv:2511.10373v2 (CC BY 4.0), distributed under the Creative Commons Attribution 4.0 International licence, as recorded in the arXiv licence metadata for this exact version and shown on the abstract page https://arxiv.org/abs/2511.10373v2. Full rights notice: licenses/arxiv-2511.10373-CC-BY-4.0.txt.\par\smallskip

\noindent\textit{Editorial scope.} Counted unit: the Thm whose source label is \texttt{thm:volume} in the article (source0.tex lines 772--778, source label \texttt{thm:volume}), delivered together with the article's own complete original proof of it (source0.tex lines 808--833, one \texttt{proof} environment of 26 lines). No mathematics is written, repaired or re-derived: statement and proof are the article's own text line for line. The proof is detached from the statement in the source: the article states the result at lines 772--778 and prints its proof at lines 808--833, and the proof environment's own header names the statement, so the association is the article's own. Required context retained uncounted: prop:translation (lines 753--760). Every definition, display and label of those lines is retained unchanged. Omitted from this package and not redistributed: 1--752, 761--771, 779--807, 834--1134. Nothing inside a retained excerpt is omitted or altered: every retained body is delivered exactly as the article's lines are written, so no omission receipt of any kind accompanies this package. Citations: the retained excerpts carry no citation command at all, so no bibliography entry of the article is transcribed and no reference is redistributed. Reference adaptation: none was needed and none was made; every reference inside the delivered unit resolves inside it, so no unresolved reference or citation is produced. Printed slips kept literal: no printed word, symbol, hypothesis or number of the counted statement or of its proof was altered. Layout and class: the article's own class is \texttt{amsart} with packages including babel, inputenc, amsmath, amsfonts, amssymb, amsthm, geometry, hyperref, cleveref, graphicx, xfrac, color, tikz, tikz-cd; the wrapper is the lane's pinned \texttt{amsart} editorial wrapper at 10pt on A4, which restates the theorem-like environments the retained text needs and adds the article's own macro definitions that the retained text needs (\texttt{\textbackslash Q}, \texttt{\textbackslash R}, \texttt{\textbackslash Z}, \texttt{\textbackslash ehr}, \texttt{\textbackslash phi}, \texttt{\textbackslash vol}), with extra wrapper packages (bm, bbm, dsfont, xspace, cancel, mathtools). The delivered numbering is the unit's own. Assets: no retained span contains a figure environment, a graphics-inclusion command, a PDF-page inclusion, a table or a TikZ picture; no image file is copied into this package, the package declares no asset dependency and no image file is redistributed with it. Licence: version arXiv:2511.10373v2 of this article is distributed under the Creative Commons Attribution 4.0 International licence, as recorded in the arXiv licence metadata for this exact version and shown on the abstract page https://arxiv.org/abs/2511.10373v2. A full rights notice is delivered at licenses/arxiv-2511.10373-CC-BY-4.0.txt. Characters: every retained excerpt is pure ASCII, so the delivered engine, XeLaTeX with its default Latin Modern text font, reports no missing character.
\begin{thm}
    \label{prop:translation}
Let $P$ be a rational polytope with denominator $q$ and let $v$ be a rational vector with denominator $q$. Then
$$\sum_{n\geq 0}|(nP-v)\cap \Z^d|t^n=\frac{h_v^*(t)}{(1-t^q)^{d+1}},$$
where $h_v^*(t)$ is a polynomial with nonnegative integer coefficients. In particular, if $P$ is a lattice polytope and $v$ an integer vector, we recover the classical Ehrhart series
$$\mathcal E(P;t) = \sum_{n\geq 0}\ehr(P;n)t^n=\frac{h^*(t)}{(1-t)^{d+1}} \, ,$$
where $h^*(t)$ is the $h^*$-polynomial of the polytope $P$.
\end{thm}

\begin{thm}\label{thm:volume}
Let $P \subseteq \R^d$ be a $d$-dimensional lattice polytope, $G$ be finite, and $\varphi:\Z^d \twoheadrightarrow G$ be surjective. As a polynomial function  in $n$, the leading coefficient of $\ehr(P,\varphi;n)$ is  \[\vol(P) \cdot \frac{1}{|G|} \sum_{g \in G} g.\] 
Equivalently, for any $g \in G$,
\[
\lim_{n \to \infty} \frac{[g]\ehr(P,\varphi;n)}{n^d}= \frac{\vol(P)}{|G|}.
\]
\end{thm}

\begin{proof}[Proof of Theorem~\ref{thm:volume}]
For simplicity of notation, define $\psi_g(n):=[g]\ehr(P,\phi;n)$ for any $g \in G$.  Fix $a_g \in \Z^d$ as any lattice point in the preimage of $g$ under $\phi$. Then we have
\[
\psi_g(n) =|\{\alpha \in nP \cap \Z^d \colon \phi(\alpha)=g\}|= |\{\alpha \in nP \cap \mathbb{Z}^d \colon \alpha \in a_g + \ker(\varphi)\}| = |nP \cap (a_g+\ker(\varphi))|.
\]
The second equality follows from the fact that $\phi(\alpha)=g$ if and only if $\alpha \in a_g+\ker(\varphi)$.

Since $\phi$ is surjective, $\ker(\phi)$ is a full-rank sublattice of $\Z^d$ of index $|G|$. Let $B$ be a $d \times d$ integer matrix whose columns form a $\Z$-basis of $\ker(\phi)$. Then $|\det(B)|=[\Z^d:\ker(\phi)]=|G|$. 
However $\alpha \in nP \cap (a_g+\ker(\phi))$ if and only if $\alpha \in nP$ and $\alpha=a_g+Bv$ for some $v \in \Z^d$. 
The invertibility of $B$ over $\Q$  gives 
\[
\psi_g(n) = |\{v \in \Z^d \colon v \in B^{-1}(nP-a_g)\}|.
\]
Let $Q$ be the full dimensional rational polytope obtained by applying the transformation $B^{-1}$ to $P$. Then we have $B^{-1}(nP-a_g)=nQ-B^{-1}a_g$ which is a translation of $nQ$ by the vector $B^{-1}a_g \in \mathbb{Q}^d$. We now have
\[
\psi_g(n)=|(nQ-B^{-1}a_g) \cap \mathbb{Z}^d|.
\]
Then it follows from Theorem~\ref{prop:translation} that $\psi_g(n)$ is a quasipolynomial in $n$ of degree $d$. Its leading coefficient is equal to 
\[
\lim_{n\rightarrow \infty} \frac{\psi_g(n)}{n^d} = \lim_{n\rightarrow \infty} \frac{1}{n^d} \left| \left(Q-\frac{1}{n}B^{-1}a_g\right) \cap \left(\frac{1}{n}\Z\right)^d\right| = \vol(Q).
\]



Finally, $\vol(Q)=\vol(B^{-1}P)=|\det(B^{-1})|\vol(P)=\dfrac{\vol(P)}{|G|}$. The result follows.
\end{proof}

\end{document}
